Distributed service-oriented applications require secure discovery, provider selection, and access control under changing resource availability. Earlier service information may not reflect a provider's suitability for a particular request. We present NDNSF, a distributed, data-driven service framework built on Named Data Networking (NDN). Its basic invocation workflow uses Request, acknowledgment (ACK), Selection, and Response Data to solicit provider participation, select executors from their feedback, and deliver results. NDNSF uses a data synchronization protocol to disseminate these messages. Provider feedback and customizable selection allow users to choose executors according to task requirements and current resource availability. To protect task requirements disclosed during discovery, NDNSF uses name-based access control with attribute-based encryption (NAC-ABE). This allows multiple authorized providers to assess a request while keeping its discovery descriptor confidential from readers lacking the required decryption capability, without enumerating provider identities. Under one ABE authority and the same set of public parameters, NDNSF also aggregates service permissions in one decryption key per entity. Separating service permissions from signing identity supports credential reuse across compatible services and applications, with baseline trust rules focused on signing authority rather than service-permission lists. We analyze credential counts and identify policy-dependent distribution and update costs, then evaluate invocation behavior against NDN Service Call (NSC) and gRPC in wired and wireless experiments. In our single-provider wired experiments with packet loss, NDNSF achieves higher completion rates than NSC and rates comparable to gRPC. With multiple intermittently available providers, it achieves higher completion rates than the fixed-provider baselines by dynamically selecting available providers.
